\chapter{Conclusions and Recommendations}

\section{Conclusions} 
This study involves several findings. The followings are the most important points and conclusions of this research:

A framework of predictive data analytics applications on the oilfield is used to state various applications. Then a project from each study is made as a proof of concept for such applications impact on the subsurface energy systems production

For real time applications,  ...

In predictive maintenance application, predictive maintenance of the electrical submersible pumps was our proposed application. In this application, sensors measurements with moving difference is applied to the dataset in order to predict pumping condition. Then a dimensionality reduction technique is used and the whole dataset points have been projected to the new lower dimensions. Finally, these new transformed data have been pipelined with supervised algorithm, which is XGBoosting in our application. The training dataset consists of inputs (PCA projected features) paired with the representative outputs (in case of ESP failure prediction the labelled outputs are the seven days before the reported failures). Each of these inputs-output pairs should be seen as a “data point” that can be used to train, validate and test the proposed model. 

On validation set,  the proposed model has a mean AUC for the 10 fold validation equal to 0.95. which in turn means that the model has an adequate performance and can be tested on the later processes against test set. 

On testing set, the proposed model can report the pre-workover and workover class with 0.8 precision and 0.6 recall. The model has high precision on testing set and hence a small number of false alarms. Of course, the relevant recall is small, which means not all the 7 days before the event are marked as yellow alarm (pre-workover and workover event). In other words, the model will report alarms with high precision but not for all days before the workover. This is acceptable because it is not necessary that all days before the event may exhibits a sign of an upcoming workover.     

For prescriptive applications, This research propose an application of an Actor-Critic reinforcement learning (RL) approach is used to optimize steam injection rates based on a problem of compositional modelling of a fluid flowing through a porous medium. The objective was to find a policy able to maximize the cumulative net present value at the end of the production horizon by optimizing the well water injections rates. The purpose of the Actor network is to estimate the value function through the mimic of the physics of the environment inputs and outputs, and to update the policy distribution in the direction suggested by the critic network. After a successive interaction between the Agent (actor critic network) and the reservoir simulation model, it is shown that the agent was able to choose a series of actions (steam injection rate versus time) aimed at maximizing the net present value of the project. The results from both, the base and the optimum policy case are compared from a physical perspective to explain the reasons leading the increase on the net present value of the whole project. 

This work represents the first step in the direction of implementing cooperative competitive multiagent actor critic reinforcement learning framework to optimize multi-pad oil wells.

\section{Recommendations and further work}

Machine learning methods involve two significant components that would affect the prediction process. The first one is the data and the second one is the developed models.


The applicability and impact on the field column summarizes the discussions 

It depicts the fitness of the application the specific conditions in the field and the aspired reward by introducing the application. Both Scales range from + “good applicability/high impact” to “+/-“ “medium applicability/ medium impact” and - “low applicability/low impact”. The applicability of Optimization of team injection is depicted as “?”, stating that further investigation would be required regarding the availability of history matched reservoir models. 


\subsection{Recommendations for forecasting problems}
This research is based on the machine learning techniques
. 
\subsection{Recommendation for the predictive maintenance}  
In this study, 

\subsection{Recommendation for the prescriptive maintenance}
